\documentclass[11pt]{article}

\usepackage{amsmath,amsfonts,amssymb}
\usepackage{graphicx}
\usepackage{booktabs}
\usepackage{hyperref}
\usepackage{geometry}
\usepackage{enumitem}
\usepackage{xcolor}
\usepackage{mdframed}
\usepackage{natbib}
\usepackage{array}
\usepackage{multirow}
\usepackage{algorithm}
\usepackage{algorithmic}
\usepackage{subcaption}

\geometry{margin=1in}

% ── Callout boxes ─────────────────────────────────────────────────────────────
\definecolor{shadeblue}{RGB}{235,243,251}
\definecolor{shadegray}{RGB}{245,245,245}
\definecolor{shadeyellow}{RGB}{255,252,230}

\mdfdefinestyle{keyresult}{
  backgroundcolor=shadeblue, linecolor=blue!40, linewidth=0.8pt,
  innerleftmargin=10pt, innerrightmargin=10pt,
  innertopmargin=8pt,  innerbottommargin=8pt,
  skipabove=8pt, skipbelow=8pt
}
\mdfdefinestyle{graybox}{
  backgroundcolor=shadegray, linecolor=gray!40, linewidth=0.8pt,
  innerleftmargin=10pt, innerrightmargin=10pt,
  innertopmargin=6pt,  innerbottommargin=6pt,
  skipabove=6pt, skipbelow=6pt
}

% ── Title ─────────────────────────────────────────────────────────────────────
\title{
  \textbf{FinStructBench:}\\[4pt]
  \large A Benchmark for Structured Information Retrieval\\
  from Financial Documents Using Graph-Verifiable Questions
}
\author{Agus Sudjianto}
\date{}

\begin{document}
\maketitle

% =============================================================================
\begin{abstract}
We introduce \textbf{FinStructBench}, a benchmark for evaluating the ability of
large language models (LLMs) to extract, aggregate, and reason over structured
information in financial documents.  Unlike existing QA benchmarks that rely on
human-annotated answers, FinStructBench generates questions \emph{automatically}
from a document's underlying knowledge graph, making ground truth
\emph{provably correct by construction}.  We define six question categories of
increasing structural complexity---exact recall, threshold checking,
cross-reference, counting, contradiction detection, and multi-hop
reasoning---and provide five benchmark instances spanning major financial
report types: model validation (SR~11-7), fair lending (ECOA/HMDA), stress
testing (CCAR/DFAST), credit portfolio review (OCC), and Basel~III capital
adequacy (Pillar~3).  Across 258 questions, a graph-based retrieval baseline
achieves 100\% accuracy by design, while Claude Sonnet~4 scores 45\%,
revealing a sharp \emph{complexity gradient}: LLM accuracy drops from 86\% on
simple threshold queries to 0\% on multi-hop reasoning.  FinStructBench is
released as an open-source toolkit enabling researchers to generate
benchmark instances from any markdown-formatted financial document.
\end{abstract}

% =============================================================================
\section{Introduction}
\label{sec:intro}

Financial regulation demands precise information extraction from structured
documents: capital ratios must be reported to basis-point accuracy, adverse
impact ratios must be computed across protected classes, and stress-test
projections must be traced across scenarios and time horizons.  Errors in
these tasks carry material consequences---regulatory findings, enforcement
actions, and capital shortfalls.

Large language models are increasingly deployed for financial document
analysis, yet existing benchmarks for financial NLP focus predominantly on
sentiment analysis~\citep{malo2014good}, named entity recognition~\citep{alvarado2015domain},
or free-form question answering~\citep{chen2021finqa, zhu2021tat} where ground
truth is established by human annotation.  These benchmarks have three
limitations in the structured-retrieval setting:

\begin{enumerate}[leftmargin=2em]
  \item \textbf{Annotation cost and error.}  Human annotators make mistakes on
    precise numeric lookups, especially in dense tables with hundreds of cells.
  \item \textbf{Limited coverage.}  Hand-crafted questions test only the
    patterns the annotator anticipated, missing systematic failure modes.
  \item \textbf{No structural baseline.}  Without a provably correct retrieval
    system, it is impossible to distinguish LLM errors from ambiguous ground truth.
\end{enumerate}

We address all three limitations with \textbf{FinStructBench}, a benchmark
construction methodology that:
\begin{itemize}[leftmargin=2em]
  \item Ingests any markdown-formatted financial document into a
    \emph{DocumentGraph} comprising exact numeric memory (ENM) entries and
    knowledge graph (KG) triples;
  \item Automatically generates questions by mining graph topology, with
    ground truth defined by graph traversal---making it provably correct;
  \item Provides a \emph{graph-based retrieval baseline} (the graph itself)
    that achieves 100\% by construction, against which any LLM or agentic
    system can be compared.
\end{itemize}

Our initial release includes 258 questions across five financial report
types and six complexity categories.  Results with Claude Sonnet~4 reveal a
monotonic complexity gradient from 86\% on threshold queries to 0\% on
multi-hop reasoning, demonstrating that current LLMs lack the structured
traversal capabilities required for reliable financial document analysis.

% =============================================================================
\section{Related Work}
\label{sec:related}

\paragraph{Financial NLP Benchmarks.}
FinQA~\citep{chen2021finqa} and TAT-QA~\citep{zhu2021tat} evaluate numerical
reasoning over financial tables, but rely on human-annotated answers and focus
on arithmetic operations (addition, percentage change) rather than structural
graph operations.  ConvFinQA~\citep{chen2022convfinqa} extends to
conversational settings but inherits the same annotation limitations.
FLUE~\citep{shah2022flue} aggregates multiple financial NLP tasks but does not
address structured retrieval from regulatory documents.

\paragraph{Table QA.}
WikiTableQuestions~\citep{pasupat2015compositional} and
HybridQA~\citep{chen2020hybridqa} benchmark table understanding but use
general-domain tables without the precision requirements of financial
regulation.  SQA~\citep{iyyer2017search} addresses sequential table queries
but does not test cross-table reasoning.

\paragraph{Document Understanding.}
DocVQA~\citep{mathew2021docvqa} and financial variants evaluate visual
document understanding, which is complementary to our text-based setting.
Recent work on long-context evaluation~\citep{li2024long} tests retrieval
from long documents but uses synthetic ``needle-in-haystack'' probes rather
than domain-specific structural queries.

\paragraph{Knowledge Graph QA.}
KGQA benchmarks~\citep{zhang2022subgraph} evaluate question answering over
pre-existing knowledge graphs.  FinStructBench differs in that the graph is
\emph{constructed from} the document, and the benchmark tests whether an LLM
can implicitly perform the same traversal from the raw text.

% =============================================================================
\section{Benchmark Construction}
\label{sec:construction}

FinStructBench's construction pipeline has three stages: document ingestion,
question generation, and evaluation.  Each stage is fully automated and
domain-agnostic.

\subsection{Document Ingestion}
\label{sec:ingestion}

Given a markdown-formatted financial document, the ingester produces a
\textbf{DocumentGraph} $\mathcal{G} = (\mathcal{E}, \mathcal{T}, \Phi)$
with three components:

\begin{description}[leftmargin=1.5em, labelindent=0em]
  \item[Exact Numeric Memory (ENM) $\mathcal{E}$:]  A key-value store
    mapping $(\textit{type}, \textit{id}) \to (\textit{value},
    \textit{hash})$ where each numeric value is stored with its SHA-256
    hash for integrity verification.  The \textit{type} is derived from
    the sanitized section heading (e.g.,
    \texttt{capital\_adequacy\_ratios}), and the \textit{id} is a
    composite key encoding the entity, any secondary label columns, and
    (for multi-column tables) the column name, separated by~\texttt{/}.
    For example, a stress test table with columns Scenario, Quarter,
    CET1, and Tier~1 produces keys such as
    (\texttt{capital\_projections},
    \texttt{Baseline/Q1\,2026/CET1}),
    ensuring unambiguous lookup even when the same metric name appears
    in multiple rows.  Each value is stored alongside its SHA-256 hash
    (computed over the IEEE-754 double representation) so that any
    silent corruption is detected at read time.

  \item[Knowledge Graph Triples $\mathcal{T}$:]  A set of
    $(\textit{head}, \textit{relation}, \textit{tail})$ triples extracted
    from table structure.  For each row the ingester emits:
    (i)~a \texttt{has\_\textit{col}} triple for every numeric cell,
    linking the entity to its string-formatted value;
    (ii)~a categorical triple for every non-entity label column;
    (iii)~a \texttt{passes} or \texttt{fails} triple if a boolean column
    is present; and
    (iv)~an \texttt{in\_section} triple linking the entity to its
    section tag.

  \item[Phase Encoders $\Phi$:]  Inequality functions that map a metric
    type to a comparison operation.  Each encoder stores the value range
    $[v_{\min}, v_{\max}]$ observed for that metric (used internally for
    geometric encoding) and exposes a \texttt{check\_inequality}$(v,
    \tau, \mathrm{op})$ method supporting $\geq$, $>$, $\leq$, and $<$.
    Encoders are inferred from column names using regex matching against a
    dictionary of 11 financial metric types (Table~\ref{tab:encoders}).
\end{description}

\begin{algorithm}[t]
\caption{Document Ingestion Pipeline}
\label{alg:ingest}
\begin{algorithmic}[1]
\REQUIRE Markdown document $D$
\ENSURE DocumentGraph $\mathcal{G} = (\mathcal{E}, \mathcal{T}, \Phi)$
\STATE Parse $D$ into sections $\{s_1, \ldots, s_k\}$ by \texttt{\#\#}/\texttt{\#\#\#} heading hierarchy
\STATE Infer phase encoders $\Phi$ from column names across all tables
\FOR{each section $s_i$}
  \FOR{each markdown table $T_j$ in $s_i$}
    \STATE Classify columns into \textit{entity}, \textit{numeric}, \textit{boolean}, \textit{label}
    \STATE Detect entity column by name heuristic; detect boolean column
    \FOR{each row $r$ in $T_j$}
      \STATE $e \gets$ entity value (or subsection name, or \texttt{row\_$i$})
      \STATE $e_{\text{full}} \gets e\,/\,\text{label}_1\,/\,\ldots\,/\,\text{label}_m$ \hfill\COMMENT{composite ID}
      \FOR{each numeric column $c$}
        \STATE $\textit{id} \gets e_{\text{full}}/c$ if multi-column, else $e_{\text{full}}$
        \STATE $\mathcal{E}[\text{section\_tag}, \textit{id}] \gets (v_c,\; \text{SHA-256}(v_c))$
        \STATE $\mathcal{T} \gets \mathcal{T} \cup \{(e, \texttt{has\_}c, v_c)\}$
      \ENDFOR
      \STATE Add label triples: $(e, \texttt{has\_}\ell, v_\ell)$ for each label column $\ell$
      \IF{boolean column value $\in$ \{true, yes, weak, fail\}}
        \STATE $\mathcal{T} \gets \mathcal{T} \cup \{(e_{\text{full}}, \texttt{fails}, \text{tag})\}$
      \ELSIF{$\in$ \{false, no, strong, pass\}}
        \STATE $\mathcal{T} \gets \mathcal{T} \cup \{(e_{\text{full}}, \texttt{passes}, \text{tag})\}$
      \ENDIF
    \ENDFOR
  \ENDFOR
  \STATE Extract bullet-list key-value pairs as metadata ENM entries and triples
\ENDFOR
\RETURN $\mathcal{G}$
\end{algorithmic}
\end{algorithm}

\paragraph{Column Classification.}
The ingester classifies each column into one of four types using
value-pattern analysis:
\begin{itemize}[leftmargin=2em]
  \item \textbf{Numeric:} columns where $>$70\% of non-empty values parse
    as floats (after stripping \$, \%, commas, and parentheses for negative
    values).
  \item \textbf{Boolean:} columns where \emph{every} non-empty value is
    one of: pass, fail, yes, no, true, false, weak, strong
    (case-insensitive).
  \item \textbf{Label:} non-numeric, non-boolean string columns.
  \item \textbf{Entity:} the primary identifier column, selected by name
    heuristic from a list of 16 common column names (e.g., ``Feature'',
    ``Segment'', ``Scenario'', ``Product''), falling back to the first
    label column.
\end{itemize}

\noindent
This classification is conservative by design: mixed columns (e.g.,
\texttt{3.2~(Pass)}) are classified as label rather than boolean or
numeric, meaning the corresponding pass/fail relationship is not
captured in the graph.  Questions are generated only from data
that \emph{is} captured, so coverage may be incomplete but ground
truth is never wrong.

\paragraph{Phase Encoders.}
Phase encoders serve two roles: (i)~they define the value range for
normalizing numeric values to $[0, 1]$, and (ii)~they provide the
threshold values and comparison operators used by the Threshold
question generator.
Table~\ref{tab:encoders} lists the 11 encoder types, their
column-name patterns, and the default threshold used for question
generation (not the encoding range).

\begin{table}[t]
\centering
\caption{Phase encoder types: column-name patterns, default question
thresholds, and comparison direction.  Thresholds are domain-standard
regulatory or statistical bounds; the Direction column shows whether
values should be above ($\geq$) or below ($\leq$) the threshold.}
\label{tab:encoders}
\small
\begin{tabular}{llcl}
\toprule
Metric Type & Pattern & Threshold & Direction \\
\midrule
AUC & \texttt{auc, auroc} & 0.5 & $\geq$ \\
Accuracy & \texttt{accuracy, precision, recall, f1} & 0.5 & $\geq$ \\
AIR & \texttt{air, impact\_ratio} & 0.8 & $\geq$ \\
Ratio & \texttt{ratio} & 0.0 & $\geq$ \\
Rate & \texttt{rate} & 0.05 & $<$ \\
Loss & \texttt{loss, charge.off} & 0.5 & $<$ \\
Capital Ratio & \texttt{capital.*ratio, cet1, tier} & 4.5 & $\geq$ \\
Importance & \texttt{importance, weight} & 0.1 & $\geq$ \\
Divergence & \texttt{divergence, psi, kl} & 0.3 & $\geq$ \\
P-value & \texttt{pvalue, p.value} & 0.05 & $<$ \\
Percentage & \texttt{percentage, pct, share} & 0.0 & $\geq$ \\
\bottomrule
\end{tabular}
\end{table}


\subsection{Question Generation}
\label{sec:question_gen}

Questions are generated automatically by mining the DocumentGraph topology.
Each generator implements a \texttt{generate(graph)} method that returns a
pool of candidate questions, each comprising: (i)~a natural-language
question string, (ii)~a deterministic graph-traversal function that
computes the ground-truth answer, (iii)~a structured LLM prompt with
format instructions, and (iv)~a typed answer format (float, int, bool,
set, or string).  We define six categories of increasing structural
complexity:

\begin{mdframed}[style=graybox]
\small
\begin{enumerate}[leftmargin=1.5em]
\item \textbf{Exact Recall} --- Look up a single ENM value by key.
  Iterates over all ENM entries and generates one question per entry.\\
  \emph{``What is the exact value of CET1 capital ratio?''}\\
  Graph operation: $\mathcal{E}[\text{capital\_ratio}, \text{CET1}].\text{value}$\\
  Answer type: \texttt{float}.  LLM is instructed to report the exact value
  with all decimal places.

\item \textbf{Threshold} --- Check whether a value meets a domain-specific
  bound.  For each ENM entry, the generator matches its category to a
  phase encoder via keyword lookup (e.g., a category containing ``capital''
  maps to the \texttt{capital\_ratio} encoder) and generates a boolean
  question for each applicable threshold.\\
  \emph{``Does the CET1 ratio meet the minimum capital requirement ($\geq$4.5\%)?''}\\
  Graph operation: $\phi_{\text{capital\_ratio}}(\mathcal{E}[\text{CET1}]) \geq 4.5$\\
  Answer type: \texttt{bool}.

\item \textbf{Cross-Reference} --- Find entities appearing in two relation
  types.  Builds an index of relation $\to$ head entities, then generates
  one question for each pair of relations whose head-entity intersection
  has between 2 and 50 elements (filtering out trivial or degenerate
  overlaps).\\
  \emph{``Which entities appear in both `risk\_rating' and `has\_npl'?''}\\
  Graph operation: $\{h : (h,r_1,\_) \in \mathcal{T}\} \cap \{h : (h,r_2,\_) \in \mathcal{T}\}$\\
  Answer type: \texttt{set\_str} (comma-separated list).

\item \textbf{Counting} --- Aggregate over triple sets.  Three sub-patterns:
  (a)~count all triples with a given relation,
  (b)~count entities with a specific (relation, tail) pair, and
  (c)~count unique head entities per relation.\\
  \emph{``How many entities have risk\_rating = Substandard?''}\\
  Graph operation: $|\{h : (h, \text{risk\_rating}, \text{Substandard}) \in \mathcal{T}\}|$\\
  Answer type: \texttt{int}.

\item \textbf{Contradiction} --- Detect pass/fail inconsistencies
  across segments.  Identifies entities that both pass and fail tests
  within the same feature group (matched by prefix on test-name
  strings).  Generates three sub-types: (a)~count of contradicting
  features, (b)~set of contradicting features, and
  (c)~per-feature boolean questions.  Only applicable to instances
  with boolean columns (model validation in our release).\\
  \emph{``Which features have both passing and failing segments?''}\\
  Graph operation: $\{f : \exists\, (\_,\texttt{passes},f) \in \mathcal{T}$
  $\land\; \exists\, (\_,\texttt{fails},f) \in \mathcal{T}\}$\\
  Answer type: \texttt{int}, \texttt{set\_str}, or \texttt{bool}.

\item \textbf{Multi-Hop} --- Chained queries combining ENM lookups.
  Two sub-patterns: (a)~\emph{single-type}: find the argmin or argmax
  entry within one ENM type and report the entity name and value;
  (b)~\emph{cross-type}: find the argmin/max in type~$t_1$, extract the
  base entity name (the first component of the composite key before
  \texttt{/}), then look up all entries for that entity in type~$t_2$.
  Cross-type questions require at least 3 entries in $t_1$ and 2 in
  $t_2$, and are only generated when the base-entity match succeeds.\\
  \emph{``Find the entity with the lowest loss rate, then look up its
  capital ratio.''}\\
  Graph operation: $k^* = \arg\min_{k} \mathcal{E}[t_1,k]$;
  then $\{(k',v) : k' \in \mathcal{E}[t_2]$, $\text{base}(k') = \text{base}(k^*)\}$\\
  Answer type: \texttt{str} (multi-line formatted).
\end{enumerate}
\end{mdframed}

\paragraph{Sampling and Validation.}
Each generator produces a pool of all structurally valid candidates from
the graph---the pool sizes range from 8 (multi-hop in Basel Capital,
which has only 27 ENM entries) to 2{,}443 (cross-reference in Credit
Portfolio, which has 2{,}334 triples).  From each pool, up to $N$
questions are sampled uniformly without replacement using a fixed random
seed (default $N=10$, seed~42) for reproducibility.  Before inclusion,
every candidate is \emph{validated}: the graph answer function is
re-executed and its output compared to the stored ground truth.
Questions where the answer is \texttt{None}, an empty set, or
inconsistent with the stored ground truth are discarded.  This
double-check guards against bugs in the answer function and ensures
that the benchmark is self-consistent.

\paragraph{Ground Truth by Construction.}
A key property of FinStructBench is that ground truth is \emph{defined} by
graph traversal, not by human annotation.  The ingester's correctness
can be verified independently: we confirm that every stored ENM value
appears verbatim in the source markdown (1{,}572/1{,}572 entries verified
across all five instances).  Given correct ingestion, the answer function
is guaranteed correct by the determinism of key lookup, pattern matching,
and set operations.  This eliminates annotation noise and enables scaling
to arbitrarily large question sets.

\subsection{Evaluation Protocol}
\label{sec:eval_protocol}

Each question is evaluated against two systems:

\begin{enumerate}[leftmargin=2em]
  \item \textbf{Graph Baseline:}  The answer function is executed
    directly on the DocumentGraph via deterministic traversal (key
    lookups, triple pattern matching, set operations).  No learning or
    training is involved---the graph achieves 100\% by construction
    and serves as the upper bound.
  \item \textbf{LLM Under Test:}  The LLM receives the full markdown
    document wrapped in \texttt{<report>} tags as context, along with a
    system prompt instructing it to answer precisely using only the
    provided data and to copy numbers exactly as they appear.
    Each question includes a format instruction specifying the expected
    answer tag (e.g., \texttt{ANSWER: 13.232}).
\end{enumerate}

\paragraph{Response Parsing.}
LLM responses are parsed using type-specific parsers with two-stage
extraction:
\begin{itemize}[leftmargin=2em]
  \item \textbf{Float:}  First attempts regex on \texttt{ANSWER:\textbackslash s*<number>};
    falls back to the last float literal in the response.
  \item \textbf{Int:}  Same two-stage strategy; fallback extracts
    integers in the range $[1, 999]$.
  \item \textbf{Bool:}  Matches \texttt{ANSWER:\textbackslash s*(true|false|yes|no)};
    falls back to keyword detection (\texttt{passes}, \texttt{does not},
    \texttt{fails}).
  \item \textbf{Set:}  Splits the \texttt{ANSWER:} line on commas and
    strips quotes.
  \item \textbf{String:}  Returns the \texttt{ANSWER:} line verbatim.
\end{itemize}

\paragraph{Scoring.}
Scoring uses exact match with the following type-specific rules:
float answers require $|a - g| < 10^{-6}$;
integer answers require equality after casting;
boolean answers require equality;
set answers require exact set equality (superset answers are also
accepted);
string answers require case-insensitive equality.
Tuple-valued answers (multi-hop) compare element-wise with float
tolerance $10^{-3}$.

% =============================================================================
\section{Benchmark Instances}
\label{sec:instances}

We release five benchmark instances covering major financial document types
encountered in US banking regulation.  Each instance is a synthetic but
structurally realistic report: we first defined the table schemas, section
structure, and entity names to mirror actual regulatory filings, then
populated numeric cells with randomly generated values drawn from
domain-appropriate distributions (e.g., capital ratios from
$\mathcal{N}(12, 3)$ clipped to $[4, 25]$, approval rates from
$\text{Beta}(5, 2)$).  All instances are generated with fixed random
seeds for reproducibility.
Table~\ref{tab:instances} summarizes their characteristics.

\begin{table}[t]
\centering
\caption{FinStructBench instance statistics.}
\label{tab:instances}
\small
\begin{tabular}{lrrrcl}
\toprule
Instance & ENM & Triples & Questions & Encoders & Regulatory Context \\
\midrule
Model Validation & 261 & 5{,}586 & 60 & 6 & SR 11-7 \\
Fair Lending & 866 & 1{,}962 & 50 & 7 & ECOA / HMDA \\
Stress Test & 207 & 1{,}355 & 50 & 4 & CCAR / DFAST \\
Credit Portfolio & 211 & 2{,}334 & 50 & 5 & OCC Guidelines \\
Basel Capital & 27 & 785 & 48 & 5 & Basel III Pillar 3 \\
\midrule
\textbf{Total} & \textbf{1{,}572} & \textbf{12{,}022} & \textbf{258} & --- & --- \\
\bottomrule
\end{tabular}
\end{table}

\paragraph{Model Validation (SR~11-7).}
A model risk management report evaluating a credit scoring model across
multiple segments.  Contains accuracy metrics (AUC, precision, recall, F1),
feature importance analysis (FANOVA), slicing performance by demographic
segments, and resolution analysis.  This is the only instance with
contradiction questions, as it contains per-feature pass/fail assessments
that can be inconsistent across segments.

\paragraph{Fair Lending (ECOA/HMDA).}
A fair lending analysis examining loan disposition patterns across protected
classes (race, gender, age).  Contains adverse impact ratios, pricing
analysis with statistical significance tests, geographic distribution
indicators, and model fairness metrics (demographic parity, equalized odds).
The largest instance by ENM entries (866), reflecting the high dimensionality
of cross-tabulated demographic data.

\paragraph{Stress Test (CCAR/DFAST).}
A capital adequacy stress test with macroeconomic scenarios (baseline,
adverse, severely adverse) projected across nine quarters.  Contains loan
loss projections by portfolio segment, pre-provision net revenue forecasts,
capital ratio trajectories, and capital action plans.

\paragraph{Credit Portfolio Review (OCC).}
A credit portfolio examination covering portfolio composition, credit quality
metrics, risk rating distributions, concentration analysis, allowance
adequacy, vintage analysis, migration matrices, and stress loss projections.

\paragraph{Basel Capital Adequacy (Pillar~3).}
A Basel~III regulatory disclosure with capital composition, risk-weighted
assets by exposure class, capital ratios, credit risk parameters (PD, LGD,
EAD), market risk (VaR), operational risk (SMA), and leverage ratios.
The smallest instance by ENM entries (27) but with complex hierarchical
table structure.


% =============================================================================
\section{Experiments}
\label{sec:experiments}

\subsection{Setup}

We evaluate Claude Sonnet~4 (\texttt{claude-sonnet-4-20250514}) against
the graph baseline across all five instances.  The LLM receives the
complete markdown document as context (500--1{,}000 lines) along with a
structured prompt for each question specifying the expected answer format.
We use 10 questions per category, sampled with seed 42.

\subsection{Main Results}

\begin{table}[t]
\centering
\caption{Overall benchmark results: Graph baseline vs.\ Claude Sonnet~4.}
\label{tab:main_results}
\small
\begin{tabular}{lcccc}
\toprule
Instance & Questions & Graph & Claude Sonnet & $\Delta$ \\
\midrule
Model Validation & 60 & 60/60 (100\%) & 26/60 (43\%) & $-57$ pp \\
Fair Lending & 50 & 50/50 (100\%) & 30/50 (60\%) & $-40$ pp \\
Stress Test & 50 & 50/50 (100\%) & 21/50 (42\%) & $-58$ pp \\
Credit Portfolio & 50 & 50/50 (100\%) & 20/50 (40\%) & $-60$ pp \\
Basel Capital & 48 & 48/48 (100\%) & 18/48 (38\%) & $-63$ pp \\
\midrule
\textbf{Total} & \textbf{258} & \textbf{258 (100\%)} & \textbf{115 (45\%)} & $\mathbf{-55}$ \textbf{pp} \\
\bottomrule
\end{tabular}
\end{table}

Table~\ref{tab:main_results} presents the overall results.  The graph
baseline achieves 100\% across all instances by construction.  Claude
Sonnet~4 achieves 45\% overall, with per-instance scores ranging from
38\% (Basel Capital) to 60\% (Fair Lending).

\subsection{Complexity Gradient}

\begin{table}[t]
\centering
\caption{Results by question category (aggregated across all instances).
Categories are ordered by decreasing LLM accuracy, revealing a monotonic
complexity gradient.}
\label{tab:by_category}
\small
\begin{tabular}{lccccc}
\toprule
Category & Graph Op. & Total & Graph & Claude & Claude \% \\
\midrule
Threshold & $\phi(v) \geq \tau$ & 50 & 50 & 43 & 86\% \\
Cross-Reference & $S_1 \cap S_2$ & 50 & 50 & 31 & 62\% \\
Exact Recall & $\mathcal{E}[k]$ & 50 & 50 & 22 & 44\% \\
Counting & $|\{h : (h,r,t)\}|$ & 50 & 50 & 17 & 34\% \\
Contradiction & $\exists\, \text{pass} \land \exists\, \text{fail}$ & 10 & 10 & 2 & 20\% \\
Multi-Hop & chain$(f_1, f_2)$ & 48 & 48 & 0 & 0\% \\
\midrule
\textbf{Total} & --- & \textbf{258} & \textbf{258} & \textbf{115} & \textbf{45\%} \\
\bottomrule
\end{tabular}
\end{table}

Table~\ref{tab:by_category} reveals a striking complexity gradient.  LLM
accuracy decreases monotonically with the structural complexity of the
required graph operation:

\begin{itemize}[leftmargin=2em]
  \item \textbf{Threshold (86\%):}  Simple inequality checks against a
    single value.  The LLM only needs to locate one number and compare it
    to a threshold---the simplest possible structured query.
  \item \textbf{Cross-Reference (62\%):}  Set intersection across two
    relation types.  Requires identifying entities in two different contexts
    and computing their overlap.
  \item \textbf{Exact Recall (44\%):}  Precise numeric lookup.  Despite
    being a single-hop operation, LLMs frequently return nearby values
    from the same table or confuse similarly-named entities.
  \item \textbf{Counting (34\%):}  Exhaustive enumeration and aggregation.
    Requires scanning all triples matching a pattern---LLMs tend to
    undercount or lose track in long tables.
  \item \textbf{Contradiction (20\%):}  Global consistency checking.
    Requires reasoning about whether a feature passes in some segments
    but fails in others---a form of negation-as-failure that LLMs
    struggle with.
  \item \textbf{Multi-Hop (0\%):}  Chained argmin/argmax followed by
    cross-type lookup.  Claude Sonnet fails on all 48 multi-hop questions
    across all five instances, demonstrating a complete inability to
    perform sequential structured traversal.
\end{itemize}

\subsection{Per-Instance Breakdown}

\begin{table}[t]
\centering
\caption{Claude Sonnet~4 accuracy by instance and category.  Each cell
shows correct/total.  Shading indicates: {\colorbox{green!15}{$\geq$80\%}},
{\colorbox{yellow!30}{40--79\%}}, {\colorbox{red!15}{$<$40\%}}.}
\label{tab:heatmap}
\small
\begin{tabular}{l*{6}{c}}
\toprule
& \rotatebox{60}{Threshold} & \rotatebox{60}{Cross-Ref} & \rotatebox{60}{Exact Recall} & \rotatebox{60}{Counting} & \rotatebox{60}{Contradiction} & \rotatebox{60}{Multi-Hop} \\
\midrule
Model Valid. & \cellcolor{green!15}8/10 & \cellcolor{yellow!30}6/10 & \cellcolor{yellow!30}7/10 & \cellcolor{red!15}3/10 & \cellcolor{red!15}2/10 & \cellcolor{red!15}0/10 \\
Fair Lending & \cellcolor{green!15}10/10 & \cellcolor{yellow!30}7/10 & \cellcolor{green!15}9/10 & \cellcolor{yellow!30}4/10 & --- & \cellcolor{red!15}0/10 \\
Stress Test & \cellcolor{green!15}10/10 & \cellcolor{green!15}8/10 & \cellcolor{red!15}0/10 & \cellcolor{red!15}3/10 & --- & \cellcolor{red!15}0/10 \\
Credit Port. & \cellcolor{green!15}9/10 & \cellcolor{yellow!30}5/10 & \cellcolor{yellow!30}5/10 & \cellcolor{red!15}1/10 & --- & \cellcolor{red!15}0/10 \\
Basel Capital & \cellcolor{yellow!30}6/10 & \cellcolor{yellow!30}5/10 & \cellcolor{red!15}1/10 & \cellcolor{yellow!30}6/10 & --- & \cellcolor{red!15}0/8 \\
\bottomrule
\end{tabular}
\end{table}

Table~\ref{tab:heatmap} shows the full instance $\times$ category breakdown.
Several patterns emerge:

\paragraph{Multi-hop is universally zero.}
Across all five instances, Claude Sonnet fails every multi-hop question
(0/48).  These questions require identifying an extremum (argmin/argmax)
within one ENM type and then looking up the corresponding entity in a
different ENM type---a two-step traversal that the LLM cannot reliably
execute from raw text.

\paragraph{Instance difficulty varies.}
Fair Lending is the easiest instance (60\%), likely because its tables
have clear demographic category labels and the AIR threshold of 0.80 is
well-known.  Basel Capital is hardest (38\%), with deeply nested tables
and entities identified by composite keys (e.g., exposure class + sub-type
+ maturity band).

\paragraph{Exact recall difficulty depends on table structure.}
Stress Test exact recall scores 0/10 despite being a single-hop operation.
This is because the stress test uses dense quarter $\times$ scenario
matrices where multiple similar values appear in adjacent cells, causing
the LLM to extract the wrong cell.  Fair Lending scores 9/10 on exact
recall because its tables have distinctive demographic labels.

\subsection{Error Analysis}
\label{sec:errors}

We analyze the 143 LLM failures across three representative error modes:

\paragraph{Adjacent Cell Confusion.}
In 37\% of exact recall failures, the LLM returns a value from a
neighboring cell in the same table.  For example, when asked for
``Q1~2026/Non-Interest Income,'' Claude returns the Non-Interest Expense
value from the same row.  This error is especially prevalent in the
stress test instance, which uses wide matrices with 9 quarterly columns.

\paragraph{Partial Enumeration.}
In 68\% of counting failures, the LLM provides a count that is
lower than the ground truth.  The model appears to enumerate entities
until it reaches a ``satisfying'' count rather than exhaustively scanning
all triples.  For example, when asked to count entities with
\texttt{risk\_rating = Substandard}, Claude finds 4 of 7 entities,
apparently stopping after processing the first table and missing
occurrences in subsequent sections.

\paragraph{Multi-Hop Decomposition Failure.}
All 48 multi-hop failures follow the same pattern: the LLM either
(a)~answers only the first hop (reports the entity with the extreme value
but not the cross-type lookup), (b)~hallucinates the second-hop value,
or (c)~confuses which entity was the extremum.  This suggests that
current LLMs lack the ability to maintain intermediate results while
executing a second query---a fundamental limitation for agentic
financial analysis.

% =============================================================================
\section{Discussion}
\label{sec:discussion}

\paragraph{Implications for Financial AI.}
The 55 percentage-point gap between graph-based retrieval and LLM
performance has direct implications for the deployment of AI in
financial regulation.  Tasks that involve precise numeric extraction,
exhaustive enumeration, or multi-step reasoning should not rely on
LLM-only approaches.  Hybrid architectures that combine LLM
understanding with structured graph retrieval are likely necessary
for production-grade financial document analysis.

\paragraph{Benchmark as a Development Tool.}
FinStructBench can serve as an evaluation harness during the development
of retrieval-augmented generation (RAG) systems, agentic frameworks, and
tool-use architectures for financial documents.  The provably correct
ground truth eliminates the need for human evaluation, enabling rapid
iteration on system design.

\paragraph{Extensibility.}
The benchmark toolkit is designed for extensibility.  New question
categories can be added by implementing a generator class with a
\texttt{generate(graph)} method.  New instances can be created from
any markdown document.  The phase encoder system can be extended with
domain-specific threshold rules.

\paragraph{Limitations.}
FinStructBench currently uses synthetic financial documents rather
than real regulatory filings.  While the synthetic reports are
designed to be structurally realistic, real documents may contain
additional complexity (e.g., footnotes, cross-references to external
regulations, redacted sections).  We plan to extend the benchmark
with anonymized real-world instances in future work.

The benchmark evaluates a single LLM (Claude Sonnet~4) in a
zero-shot, full-context setting.  Future work should evaluate
multiple models, few-shot prompting, chain-of-thought reasoning,
and tool-augmented architectures.

% =============================================================================
\section{Released Artifacts}
\label{sec:artifacts}

FinStructBench is released as an open-source Python package with
the following components:

\begin{enumerate}[leftmargin=2em]
  \item \textbf{Toolkit} (\texttt{finstructbench/}): Document ingester,
    six question generators, scoring module, LLM caller, and benchmark
    orchestrator.
  \item \textbf{Instances} (\texttt{instances/}): Five markdown financial
    reports with reproducible generator scripts.
  \item \textbf{Results} (\texttt{results/}): Baseline results with
    per-question details in JSON format.
  \item \textbf{CLI}: \texttt{python -m finstructbench run <doc.md>}
    with options for category limit, model, and graph-only mode.
\end{enumerate}

% =============================================================================
\section{Conclusion}
\label{sec:conclusion}

We presented FinStructBench, a benchmark for structured information
retrieval from financial documents with provably correct ground truth.
Our results reveal a sharp complexity gradient in LLM performance:
Claude Sonnet~4 achieves 86\% on simple threshold checks but 0\% on
multi-hop reasoning, with an overall score of 45\% against the 100\%
graph baseline.  These findings demonstrate that current LLMs are
insufficient for reliable financial document analysis and motivate
the development of hybrid graph-augmented architectures.  FinStructBench
provides a rigorous, extensible evaluation framework for measuring
progress toward this goal.

% =============================================================================
\bibliographystyle{plainnat}
\begin{thebibliography}{20}

\bibitem[Alvarado et~al.(2015)]{alvarado2015domain}
Alvarado, J.~C.~S., Verspoor, K., and Baldwin, T. (2015).
\newblock Domain adaption of named entity recognition to support credit risk
  assessment.
\newblock In \emph{Proceedings of the Australasian Language Technology
  Association Workshop}, pp.~84--90.

\bibitem[Chen et~al.(2020)]{chen2020hybridqa}
Chen, W., Zha, H., Chen, Z., Xiong, W., Wang, H., and Wang, W.~Y. (2020).
\newblock {HybridQA}: A dataset of multi-hop question answering over tabular
  and textual data.
\newblock In \emph{Findings of EMNLP}, pp.~1026--1036.

\bibitem[Chen et~al.(2021)]{chen2021finqa}
Chen, Z., Chen, W., Smiley, C., Shah, S., Borber, I., Leu, S., Zhu, Y.,
  Kumar, S., Zhang, J., and Wang, W.~Y. (2021).
\newblock {FinQA}: A dataset of numerical reasoning over financial data.
\newblock In \emph{Proceedings of EMNLP}, pp.~7862--7872.

\bibitem[Chen et~al.(2022)]{chen2022convfinqa}
Chen, Z., Li, S., Smiley, C., Ma, Z., Shah, S., and Wang, W.~Y. (2022).
\newblock {ConvFinQA}: Exploring the chain of numerical reasoning in
  conversational finance question answering.
\newblock In \emph{Proceedings of EMNLP}, pp.~6279--6292.

\bibitem[Iyyer et~al.(2017)]{iyyer2017search}
Iyyer, M., Yih, W., and Chang, M. (2017).
\newblock Search-based neural structured learning for sequential question
  answering.
\newblock In \emph{Proceedings of ACL}, pp.~1821--1831.

\bibitem[Li et~al.(2024)]{li2024long}
Li, Y., Packer, C., and Gonzalez, J.~E. (2024).
\newblock How long can open-source {LLMs} truly promise on context length?
\newblock \emph{arXiv preprint arXiv:2402.17688}.

\bibitem[Malo et~al.(2014)]{malo2014good}
Malo, P., Sinha, A., Korhonen, P., Wallenius, J., and Takala, P. (2014).
\newblock Good debt or bad debt: Detecting semantic orientations in economic
  texts.
\newblock \emph{Journal of the Association for Information Science and
  Technology}, 65(4):782--796.

\bibitem[Mathew et~al.(2021)]{mathew2021docvqa}
Mathew, M., Karatzas, D., and Jawahar, C.~V. (2021).
\newblock {DocVQA}: A dataset for {VQA} on document images.
\newblock In \emph{Proceedings of WACV}, pp.~2200--2209.

\bibitem[Pasupat \& Liang(2015)]{pasupat2015compositional}
Pasupat, P. and Liang, P. (2015).
\newblock Compositional semantic parsing on semi-structured tables.
\newblock In \emph{Proceedings of ACL}, pp.~1470--1480.

\bibitem[Shah et~al.(2022)]{shah2022flue}
Shah, R., Chawla, K., Eidnani, D., Shah, A., Du, W., Chava, S., Raman, N.,
  Smiley, C., Chen, J., and Yang, D. (2022).
\newblock When {FLUE} meets {FLANG}: Benchmarks and large pre-trained language
  model for financial domain.
\newblock In \emph{Proceedings of EMNLP}, pp.~2322--2335.

\bibitem[Zhang et~al.(2022)]{zhang2022subgraph}
Zhang, J., Chen, B., Zhang, L., Ke, X., and Ding, H. (2022).
\newblock Neural, symbolic and neural-symbolic reasoning on knowledge graphs.
\newblock \emph{AI Open}, 3:109--124.

\bibitem[Zhu et~al.(2021)]{zhu2021tat}
Zhu, F., Lei, W., Huang, Y., Wang, C., Zhang, S., Lv, J., Feng, F., and
  Chua, T. (2021).
\newblock {TAT-QA}: A question answering benchmark on a hybrid of tabular and
  textual content in finance.
\newblock In \emph{Proceedings of ACL}, pp.~3277--3287.

\end{thebibliography}

\end{document}
